\section{Appendix: Scope and Qualifications}

This study concerns the ten propers of the Nineteenth Sunday after
Pentecost, class II, in the 1962 Roman Missal, printed pp.~411--412,
for the general-calendar occurrence on 4 October 2026. It includes
no local calendar, institute proper, commemoration of St Francis,
or optional external solemnity of the Rosary. It does not determine
current permission to use the book in a particular church. The
Ordinary is not reproduced; the Creed and Trinity Preface are noted
as appointments.

The normalized Latin was collated against the target facsimile and
the public-domain Pustet 1862 antecedent. The target edition itself
is not represented as public domain. Cummiskey's 1861 historical
English supplies the Introit antiphon and the three orations;
Douay--Rheims--Challoner supplies biblical English. Liturgical
introductions, verbal adaptations, and abbreviated conclusions
remain distinguishable. No new English prayer or translation has
been supplied. These texts aid study and are not an approved
English altar-book edition.

The reception study uses the checked loci of Augustine, Gregory,
Chrysostom, Bellarmine, Aquinas, and Schuster. Their individual
arguments are attributed to their works; the two whole-formulary
readings are editorial syntheses. Schuster's matching Sunday
commentary is direct reception of the assembled Mass. Ancient
compilation history is not established here. The reception sweep
is bounded, not exhaustive; historical translations and inspected
Latin transcriptions are not fresh critical editions. Aquinas's
inspected OCR supports paraphrase only. The evidence and textual
limits are recorded in \texttt{propers/verified.md},
\texttt{research/scope.md}, and \texttt{research/interpretations.md}.

The historical appendix preserves the available chronology's
alternatives and unresolved Gospel event date. It does not give
the broad Psalter boundary the force of a date for an individual
poem. The study is an AI-assisted Catholic study aid, not an
ecclesiastical judgment or an independently issued critical edition.

\section{References}

\begin{description}[style=nextline,leftmargin=0pt,labelsep=0pt,itemsep=.45em]
\item[Liturgical text and historical English]
\work{Missale Romanum}, Vatican editio typica, 1962, pp.~411--412,
nn.~1679--1688; general rubrics 91, 95, 111b and Mass rubrics 358b,
434b for the occurrence boundaries. CMAA facsimile,
\url{https://media.churchmusicassociation.org/pdf/missale62.pdf}.
\work{Missale Romanum}, Pustet, Ratisbon, 1862, pp.~351--352,
underlying Latin antecedent.
\work{Roman Missal Translated into the English Language for the Use
of the Laity}, Cummiskey, Philadelphia, 1861, ``Nineteenth Sunday
after Pentecost,'' Introit and three orations, pp.~445--448;
\url{https://archive.org/details/romanmissaltran00churgoog}.

\item[Scripture]
Douay--Rheims--Challoner, Project Gutenberg eBook 1581,
Psalms 77, 104, 118, 137, 140; Matthew 21--22; Ephesians 4--5.
The appointed verses are identified at each text. Psalm numbering
is Vulgate first, common modern numbering in parentheses.
\url{https://www.gutenberg.org/ebooks/1581}.

\item[Augustine]
Sermon 90, \S\S1--10; English Sermon XL, Benedictine XC, in
\work{Nicene and Post-Nicene Fathers}, first series, vol.~6.
\work{Enarrationes in Psalmos} 77.1--3; 104.1--3; 118, opening
Aleph exposition, \S\S1--8; 137.10--13; 140.1--7, in NPNF I.8.
The English headings are respectively Psalms 78, 105, 119, 138,
and 141. Historical English inspected in the CCEL text witnesses.
\url{https://www.ccel.org/ccel/schaff/npnf106.html};
\url{https://www.ccel.org/ccel/schaff/npnf108.html}.

\item[Gregory the Great]
\work{Homiliae in Evangelia} II.38, especially \S\S1--3, 7,
9--14. Historical Latin, Wikisource transcription; the heading's
pericope label stops at verse 13, but the homily prints and expounds
the concluding called/chosen saying.
\url{https://la.wikisource.org/wiki/Homiliae_in_Evangelia}.

\item[John Chrysostom]
\work{Homilies on Ephesians} XIII, on 4:22--24, and XIV, on
4:25--28, in NPNF I.13. Historical English, CCEL witness.
\url{https://www.ccel.org/ccel/schaff/npnf113.html}.

\item[Robert Bellarmine]
\work{A Commentary on the Book of Psalms}, translated by John
O'Sullivan, 1866, Psalms 77:1--8; 104:1--2; 118:1--8;
137:7--8; 140:1--5. Historical English transcription, used in
paraphrase, \url{https://www.ecatholic2000.com/bellarmine/commentary.shtml}.

\item[Thomas Aquinas]
\work{In omnes S. Pauli Apostoli epistolas commentaria}, Dessain,
1857, vol.~II, \work{Super Ephesios} 4, lectio 7, p.~329.
Inspected Latin OCR; the bounded account of renewal is paraphrased.

\item[Blessed Ildefonso Schuster]
\work{The Sacramentary}, English edition, Burns Oates \& Washbourne,
1927, vol.~III, ``Nineteenth Sunday after Pentecost,'' pp.~171--174.
The matching formulary's theological exposition is used; its
tentative historical reconstruction does not establish the
appointments' origins.

\item[Scriptural chronology and location]
Paulin Ladeuze, ``Epistle to the Ephesians,'' \work{Catholic
Encyclopedia} 5 (1909), ``Date and place of composition; occasion''
and destination discussion;
\url{https://www.newadvent.org/cathen/05485a.htm}.
Ferdinand Prat, ``St. Paul,'' vol.~11 (1911), ``Chronology,''
concluding table; \url{https://www.newadvent.org/cathen/11567b.htm}.
E. Jacquier, ``Gospel of St. Matthew,'' vol.~10 (1911), ``Date
and place of composition'';
\url{https://www.newadvent.org/cathen/10057a.htm}.
Alfred Durand, ``The Synoptics,'' vol.~14 (1912), ``Origin,''
Aramaic original and Greek rendering;
\url{https://www.newadvent.org/cathen/14530a.htm}.
USCCB, NABRE introductions to Psalms and Matthew, consulted
28 September 2026: \url{https://bible.usccb.org/bible/psalms/0}
and \url{https://bible.usccb.org/bible/matthew/0}.
The introductions supply bounded historical judgments, not
liturgical English.
\end{description}
